\textbf{1.} Le point $C$ appartient au cercle de diamètre $[AB]$ :
d'après le théorème de l'angle inscrit (cas de l'angle droit),
$\widehat{ACB} = 90^\circ$, donc $ABC$ est rectangle en $C$.

\medskip

\textbf{2.} Dans le triangle $ACH$, rectangle en $H$, on a
$\widehat{HAC} = 90^\circ - \widehat{ACH}$. Par ailleurs
$\widehat{HCB} = 90^\circ - \widehat{ACH}$ puisque
$\widehat{ACB} = 90^\circ$. Donc $\widehat{HAC} = \widehat{HCB}$, et
les triangles rectangles $AHC$ et $CHB$ sont semblables. Il en résulte
\[
\frac{AH}{CH} = \frac{CH}{HB},
\qquad\text{c'est-à-dire}\qquad
CH^2 = AH \cdot HB.
\]

\medskip

\textbf{3.} On reporte sur une droite deux segments consécutifs
$[AH]$ et $[HB]$ de longueurs $x$ et $y$. On trace le cercle de
diamètre $[AB]$, puis la perpendiculaire à $(AB)$ en $H$ : elle coupe
le cercle en un point $C$, et d'après la question précédente
$CH = \sqrt{xy}$. \hfill $\blacksquare$
